\documentclass[11pt]{article}
\usepackage[margin=1in]{geometry}
\usepackage{amsmath,amssymb,booktabs,graphicx,url}
\usepackage[colorlinks=true,linkcolor=blue,citecolor=blue,urlcolor=blue]{hyperref}
\usepackage{times}

\title{\textbf{TIM: Adapter Placement by\\Task Interference Mapping}}
\author{Alberto Acedo\\ \small Biome Makers Inc.}
\date{August 2026}

\begin{document}
\maketitle

\begin{abstract}
\noindent Low-rank adapters are inserted into the query and value projections of
every attention layer. That is the field's default, and it comes from an
ablation in the original LoRA paper rather than from anything measured about the
tasks at hand. We give a method, TIM (Task Interference Mapping), for choosing the placement
instead of inheriting it. Before any training, TIM passes a few batches of each
task through the model, takes principal angles between the two gradient
subspaces per module, and adapts the modules where the tasks overlap least. On Llama-3-8B fine-tuned from code to
prose, this retains $36.8\%$ more of the model's original capability than the
default placement with the same number of adapted modules (9 of 10 seeds,
Wilcoxon $p = 0.006$), with plain LoRA in both arms. The diagnostic takes minutes
and trains nothing. Its ordering places \texttt{o\_proj} first and
\texttt{q\_proj} second across three model families and three task pairs, while
the last two positions do not transfer, which is why we present it as a
measurement to run rather than a rule to adopt.
\end{abstract}

\section{The problem: a default that does not know your tasks}

Fine-tuning a language model on a new task degrades what it already did well.
Low-rank adaptation~\cite{hu2022} limits the damage by training a small number
of added parameters instead of all of them, and it has become the standard way
to specialise a model.

Where those parameters go is a choice. The original LoRA paper ablates over
which projections to adapt, on GPT-3 and at a fixed parameter budget, and
reports that query and value together perform best; the practice has been
inherited since and the major implementations default to it.

That choice has been revisited, and by several groups: Section~\ref{sec:sota}
surveys them. What none of them use, as far as we can tell, is the quantity that
governs how much a model forgets, which is how much the \emph{two specific tasks
at hand} overlap. Existing placement criteria score a module by its own
importance, by how much its features would move, or by alignment with the
pre-trained weights --- all properties of the model and the new task, not of the
relationship between the old task and the new one.

This paper asks whether selecting on that relationship works, and whether it can
be measured cheaply enough to be worth measuring.

\section{The theory it rests on}

A recent result makes the question answerable. Steele~\cite{steele2026} reports
that forgetting under LoRA follows
\begin{equation}
\mathcal{F} = \alpha\left(1 - \cos^2\theta_{\min}\right) + \beta,
\label{eq:steele}
\end{equation}
where $\theta_{\min}$ is the minimum principal angle between the gradient
subspaces of the two tasks, validated on Split-CIFAR100 with ViT-LoRA and on
sequential GLUE with RoBERTa-LoRA. The relationship is descriptive: how much a
model forgets is governed by how much the two tasks share.

Two things follow that the original work does not pursue. The governing quantity
is a property of the \emph{pair of tasks and the module}, so it varies across
modules within one model. And it can be computed \emph{before} training, because
gradients require only a forward and a backward pass and no parameter updates.

Together those give a procedure: measure the overlap per module, adapt where it
is smallest. Whether that procedure works is a separate question from whether
Equation~(\ref{eq:steele}) holds, and it is the one we answer.

\section{The method}

TIM takes a model and two tasks and returns a per-module map of how much they
interfere. Given those, we accumulate the gradient of every candidate
weight matrix over $n$ batches of each and take the $r$-dimensional row space of
each gradient by singular value decomposition, with $r$ the LoRA rank that would
be used. Two choices in how that comparison is made are not cosmetic, and we
state them because the measurement does not work without them.

\paragraph{The statistic is the dimension of overlap, not the leading cosine.}
The leading principal cosine corresponds to $\theta_{\min}$ and is what
Equation~(\ref{eq:steele}) uses, but it saturates at one as soon as two
subspaces share a single direction: it cannot distinguish sharing one direction
from sharing eight. We use $\sum_i \cos^2\theta_i$, the effective dimension of
the overlap. On constructed subspaces sharing $k$ of $8$ directions, the leading
cosine returns $1.0000$ for every $k \geq 1$ while this sum returns $1.03$,
$2.01$, $4.01$, $6.00$ and $8.00$ for $k = 1, 2, 4, 6, 8$.

\paragraph{The reference is a same-task ceiling, not a random-subspace null.}
Two gradients taken from the same model share directions dictated by that
model's structure regardless of which task produced them, so comparing against
random subspaces asks a question whose answer is trivially yes. We split the
first task in half and measure the overlap between halves: that is the most two
gradients can share with no task difference between them. What informs is the
drop from that ceiling. The ceiling itself varies substantially across modules,
and since a module with a low ceiling cannot drop far, we rank by the drop
\emph{relative} to the ceiling.

\paragraph{Scope of the candidate set.} We rank the four attention
projections. Nothing in the criterion is specific to them, and other placement
work reports MLP modules mattering~\cite{hayou2025,domlora2026}, so extending
the candidate set is straightforward; we have not measured that case and do not
claim anything about it.

The whole measurement is a forward and backward pass over a few batches, with no
optimiser and no parameter updates. It takes minutes on one GPU.

\section{Does it work?}

We fine-tune Llama-3-8B from code to prose and measure retained code capability
by HumanEval pass@1 over ten random seeds; the protocol is in
Appendix~\ref{app:proto}. The placement TIM suggests is the quartile of largest relative drop: 17 \texttt{o\_proj}, 10 \texttt{q\_proj} and 4 \texttt{k\_proj}
modules, and no \texttt{v\_proj} at all --- already at odds with the default,
since \texttt{v\_proj} is one of the two the default adapts. The comparison arm
uses query and value projections across layers spread through the model, cut to
the \emph{same count} of 31 modules, so that what varies between arms is
location and not capacity.

\begin{table}[ht]
\centering
\caption{Ten seeds, same session, same hardware, plain LoRA in both arms with no
penalty of any kind. \emph{Absolute} is HumanEval pass@1 after the new task; the
retention ratio divides it by the value after the first task. We report both
because the ratio is sensitive to variation in its denominator.}
\label{tab:main}
\begin{tabular}{lcccc}
\toprule
Placement & Modules & Retention & Absolute & vs.\ default \\
\midrule
Default (\texttt{q},\texttt{v}) & 31 & 54.1\% & 0.151 & --- \\
\textbf{Measured} & 31 & \textbf{75.6\%} & \textbf{0.206} &
  \textbf{9/10}, sign $p{=}0.011$, Wilcoxon $p{=}0.006$ \\
\bottomrule
\end{tabular}
\end{table}

The gain is $+36.8\%$ relative in absolute capability. Both arms end the first
task at the same capability ($0.274$ and $0.285$), so the retention ratio is not
inflated by a ceiling effect; the small difference runs against the measured arm
if anything.

\paragraph{The ordering transfers; its tail does not.} Running the diagnostic
alone --- no training --- on Llama-3-8B with two further task pairs, and on
Mistral-7B and Qwen2.5-7B with code to prose, puts \texttt{o\_proj} first and
\texttt{q\_proj} second in all five combinations.

\begin{table}[ht]
\centering
\caption{Relative drop from the same-task ceiling, averaged per module type.
Diagnostic only: no training was performed for these rows.}
\begin{tabular}{lcccc}
\toprule
Model and task pair & \texttt{o\_proj} & \texttt{q\_proj} & \texttt{k\_proj} & \texttt{v\_proj} \\
\midrule
Llama-3-8B, code $\to$ prose & 0.751 & 0.683 & 0.648 & 0.584 \\
Llama-3-8B, code $\to$ math & 0.698 & 0.681 & 0.646 & 0.566 \\
Llama-3-8B, prose $\to$ math & 0.511 & 0.403 & 0.213 & 0.244 \\
Mistral-7B, code $\to$ prose & 0.717 & 0.395 & 0.394 & 0.336 \\
Qwen2.5-7B, code $\to$ prose & 0.592 & 0.445 & 0.322 & 0.367 \\
\bottomrule
\end{tabular}
\end{table}

The last two positions swap in two of the five, and the size of the separation
varies widely. Across the three model families on code to prose, the gap between
the first and second module type runs from $0.068$ on Llama to $0.32$ on
Mistral, a factor of about five; across all five combinations it falls as low as
$0.017$, on Llama with code to math. The head of the ordering is not an artefact
of one model or one task pair; the tail is not stable and the magnitude is not
transferable at all. That is why the
diagnostic is something to run on a given model and task pair rather than a
fixed recommendation to adopt.

\section{Does it combine with a weight-space regulariser?}
\label{sec:combina}

A separate line of our own work penalises a topological property of the weight
graph during training and improves retention on the same
protocol~\cite{acedo2026omegas}. The two act at different points --- one chooses
where trainable parameters go, the other constrains how they end up --- so we
asked whether they compose, and measured it rather than assuming.

The answer is that we cannot tell. Adding the penalty on top of the TIM placement gives 6 of 10 seeds on the retention ratio and 7 of 10 on absolute
capability, with mean differences of $+0.047$ and $+0.027$. Neither approaches
significance ($p = 0.85$ and $p = 0.28$) and both sit below the run-to-run
standard deviation of $0.104$ reported in Section~\ref{sec:noise}. The direction
is positive on both measures, and that is the whole of what the data support.

One detail of that experiment generalises beyond it. We first ran the comparison
at the penalty strength selected for a configuration with twice as many adapted
modules, and the result came out \emph{negative}. A sweep on held-out seeds
showed that strength to be ten times too strong for half the adapted modules,
and the sign reversed on recalibration. Any penalty strength chosen for one
configuration should be re-derived for another, and a negative result obtained
without doing so says nothing.

\section{What this does not establish}

\paragraph{A second training validation.} We attempted one on Mistral-7B and
could not complete it: that base model's HumanEval capability is near zero, so
retention has no defined denominator. A second validation needs a model with
real capability on the retained task and remains open. Steele's law was
validated on other architectures, but it is a law about forgetting, not about
this selection procedure, so it does not stand in for that.

\paragraph{How much variation this setting carries.}
\label{sec:noise}
Repeating an identical configuration --- same seed, same code, same hardware ---
gives a standard deviation of $0.104$ in retention ratio, with individual repeats
as far apart as $0.596$ and $0.793$. The seed does not identify the run, so a
seed-paired test compares runs differing by more than the seed, and any mean
difference below roughly $0.066$ is not distinguishable at ten seeds. That GPU
non-determinism dominates seed variance is documented in vision~\cite{morin2020}
and reinforcement learning~\cite{nagarajan2018}; we have not found it quantified
for low-rank fine-tuning of language models, and it bounds every seed-paired
comparison in this literature, this paper included.

\section{Where this sits}
\label{sec:sota}

Adapter placement is an active question and several groups have attacked it.
Table~\ref{tab:sota} separates them by what they measure, when they measure it,
and what they optimise, because those axes are what distinguish the approaches
and because on two of them this work is alone.

\begin{table}[ht]
\centering
\small
\caption{Methods that decide where or how to adapt. \emph{Signal} is the
quantity scored, \emph{timing} is when it is computed, \emph{objective} is what
the method is optimised for. Every entry scores a property of the model or of
the new task; none scores the \emph{relationship} between the retained task and
the new one, which is what Equation~(\ref{eq:steele}) identifies as governing
forgetting.}
\label{tab:sota}
\begin{tabular}{llll}
\toprule
Method & Signal & Timing & Objective \\
\midrule
LoRA ablation~\cite{hu2022} & downstream accuracy & once, offline & new-task accuracy \\
AdaLoRA~\cite{zhang2023} & adapter singular values & during training & new-task accuracy \\
PLoP~\cite{hayou2025} & normalised feature norm & before training & new-task accuracy \\
FLoE~\cite{floe2025} & Fisher information & before training & new-task accuracy \\
Aletheia~\cite{aletheia2026} & gradient magnitude & before training & new-task accuracy \\
DomLoRA~\cite{domlora2026} & dominant module & before training & new-task accuracy \\
Amazon Nova~\cite{amazon2026} & benchmark sweep & once, offline & accuracy / latency \\
\midrule
O-LoRA~\cite{wang2023} & past-task subspace & during training & retention \\
OPLoRA~\cite{oplora2025} & top-$k$ singular dirs.\ of $W_0$ & during training & retention \\
Omega-S~\cite{acedo2026omegas} & weight-graph topology & during training & retention \\
\midrule
\textbf{TIM (this work)} & \textbf{task-pair gradient overlap} & \textbf{before training} & \textbf{retention} \\
\bottomrule
\end{tabular}
\end{table}

\paragraph{Placement methods optimise for the new task, not for what is kept.}
The upper block chooses where to adapt, and every one scores for downstream
accuracy. PLoP~\cite{hayou2025} identifies module types by how their feature
norms align with the pre-trained representation, arguing that the least aligned
contribute most to adaptation. FLoE~\cite{floe2025} selects layers by Fisher
information, Aletheia by gradient magnitude~\cite{aletheia2026}, DomLoRA by
reducing to a single dominant module~\cite{domlora2026}. These answer a
different question: given that I want to learn a new task well and cheaply,
where should the capacity go? A module that is ideal for learning fast need not
be one where learning does least damage, and none of these works reports
retention.

The largest published sweep we know of is Amazon's on their Nova 2.0 Lite
model~\cite{amazon2026}, and it is worth reporting because it agrees with us by
another route: they find \texttt{o\_proj} alone never fails outright on any
task, typically lands within a few points of the best configuration, and is an
attractive default. Our measurement fills 17 of its 31 slots with
\texttt{o\_proj} modules, arriving there from task overlap rather than from a
benchmark sweep.

\paragraph{Retention methods act during training, not on placement.} The middle
block optimises for what is kept, which is our objective, but all of them
intervene while training rather than before it. O-LoRA~\cite{wang2023}
constrains updates to be orthogonal to past-task subspaces.
OPLoRA~\cite{oplora2025} decomposes the frozen weights by SVD and projects
updates into the orthogonal complement of the top-$k$ singular subspace, proving
that this preserves the top-$k$ singular triples exactly, and introduces a
metric $\rho_k$ for how much an update aligns with dominant directions.

\paragraph{OPLoRA in particular is complementary, not competing.} It is the
closest work to ours in problem and in tooling --- both reason about subspaces,
both target forgetting --- so the distinction is worth stating precisely. OPLoRA
protects the dominant singular directions \emph{of the pre-trained weights}, a
quantity that does not depend on which tasks are involved; we measure where the
\emph{two tasks' gradients} overlap, which does. And the two act at different
points: OPLoRA modifies the update inside modules that have already been chosen,
whereas this work chooses which modules to instrument at all. Doing both --- select placement with TIM, then project updates within the
selected modules ---
is coherent and untested, and we would be glad to see it tried.

\paragraph{Why we report no head-to-head numbers.} We compare accuracy against
none of the above, because as the evaluations currently stand the comparison
would not mean anything. OPLoRA fine-tunes LLaMA-2 and Qwen2.5 on MetaMathQA and
measures forgetting on ARC and SIQA; PLoP and FLoE report downstream accuracy on
their own suites; we fine-tune Llama-3-8B from code to prose and measure
HumanEval retention. Putting those numbers in one table would be putting
different quantities, measured on different models, in one column. A real
comparison means re-running the methods under one protocol. We have not done it,
and it is the thing we would most like someone to do.

\paragraph{How much orthogonalisation has to offer here.}
Steele~\cite{steele2026} observes that orthogonalisation-based methods add
little where natural orthogonality is already high. Our measurement finds
exactly that in seven of eight principal directions, which suggests the room for
projection methods varies by setting as much as the room for placement does.

\section{Conclusion}

The default placement for low-rank adapters is a default: selected once, on one
model, and carried since. TIM replaces it with a measurement that takes minutes, and in the case we
tested doing so retains a third more capability. The
ordering that measurement produces is stable at its head across three model
families and three task pairs and unstable at its tail, which makes it a
diagnostic to run rather than a rule to copy, and that is how we offer it.

\section*{Patent}

The method described here is covered by USPTO patent application
No.~64/121,656, filed 29 July 2026. The released code is under AGPL-3.0 for
non-commercial research and academic use, with evaluation inside a commercial
organisation granted explicitly; commercial use is licensed separately.

\section*{Availability}

Code, per-seed data and a reporting template are at\\
\url{https://github.com/BiomeMakers/TIM-OmegaS}. \textbf{We would like
others to run the diagnostic on models we have not tried}, and a report that the
ordering comes out differently is as useful to us as one that confirms it.

\appendix
\section{Protocol}
\label{app:proto}

Llama-3-8B, LoRA rank 8, $\alpha = 16$, dropout $0.1$. Task A is Python from
CodeSearchNet, task B is OpenWebText. Capability is HumanEval pass@1, measured
after task A and again after task B. Ten seeds: 42, 123, 456, 789, 1011, 2022,
3033, 4044, 5055, 6066. The diagnostic uses 32 batches per task at sequence
length 256, split into halves of 16 for the same-task ceiling, and subsamples
1024 rows per module. All arms in Table~\ref{tab:main} were measured in one
session on the same hardware; per-seed values are in
\texttt{results/placement.json}.

\begin{thebibliography}{9}

\bibitem{hu2022}
Hu E J, Shen Y, Wallis P, Allen-Zhu Z, Li Y, Wang S, Wang L, Chen W (2022).
LoRA: Low-Rank Adaptation of Large Language Models.
\emph{ICLR}.

\bibitem{steele2026}
Steele B (2026).
Subspace Geometry Governs Catastrophic Forgetting in Low-Rank Adaptation.
\emph{arXiv:2603.02224}.

\bibitem{acedo2026omegas}
Acedo A (2026).
Omega-S: A Functional Resilience Index for LLM Fine-Tuning.
\emph{arXiv:2608.03887}.

\bibitem{zhang2023}
Zhang Q, Chen M, Bukharin A, He P, Cheng Y, Chen W, Zhao T (2023).
AdaLoRA: Adaptive Budget Allocation for Parameter-Efficient Fine-Tuning.
\emph{ICLR}.

\bibitem{wang2023}
Wang X, Chen T, Ge Q, Xia H, Bao R, Zheng R, Zhang Q, Gui T, Huang X (2023).
Orthogonal Subspace Learning for Language Model Continual Learning.
\emph{EMNLP Findings}.

\bibitem{hayou2025}
Hayou S, Ghosh N, Yu B (2025).
PLoP: Precise LoRA Placement for Efficient Finetuning of Large Models.
\emph{arXiv:2506.20629}.

\bibitem{oplora2025}
Xiong Y, Xie X (2025).
OPLoRA: Orthogonal Projection LoRA Prevents Catastrophic Forgetting during
Parameter-Efficient Fine-Tuning.
\emph{arXiv:2510.13003}.

\bibitem{floe2025}
Wang X, Gao L, Wang H, Zhang Y, Zhao J (2025).
FLoE: Fisher-Based Layer Selection for Efficient Sparse Adaptation of Low-Rank
Experts.
\emph{arXiv:2506.00495}.

\bibitem{aletheia2026}
Saket A (2026).
Aletheia: Gradient-Guided Layer Selection for Efficient LoRA Fine-Tuning Across
Architectures.
\emph{arXiv:2604.15351}.

\bibitem{domlora2026}
Zhang S, He R, Fang D, Tan X, Chen K, Zhuang H (2026).
Rethinking Adapter Placement: A Dominant Adaptation Module Perspective.
\emph{arXiv:2605.06183}.

\bibitem{amazon2026}
Amazon Science (2026).
Optimizing LoRA target module selection for efficient fine tuning.
\emph{Amazon Science Blog}.
\url{https://www.amazon.science/blog/optimizing-lora-target-module-selection-for-efficient-fine-tuning}

\bibitem{morin2020}
Morin M, Willetts M (2020).
Non-determinism in TensorFlow ResNets.
\emph{arXiv:2001.11396}.

\bibitem{nagarajan2018}
Nagarajan P, Warnell G, Stone P (2018).
Deterministic Implementations for Reproducibility in Deep Reinforcement
Learning.
\emph{arXiv:1809.05676}.

\end{thebibliography}
\end{document}
